\begin{lemma}[Existence of the prescribed parallel-extension color]
Let \(\mathcal T\) be a certified Tashkinov series with finite data
\(\mathcal D\), let \(i\) be one of its coloring indices, and let
\[
 Q=(\pi,v,f,u,\delta)
\]
be an admissible maximum-defective choice for
\[
 \bigl(T_{i+1},D_i,\varphi_i\bigr).
\]
If \(\Delta(G)+1\leq k\), then there is a color \(\gamma\) satisfying
the history-sensitive predicate
\[
                 \operatorname{ParallelExtensionColor}
                    (\mathcal D,i,\pi,v,\gamma).
\]
Thus, if a positive PE slot at most \(i\) has support \(v\), \(\gamma\)
is the delta color of the latest such slot; if no such slot exists,
\(\gamma\) is an arbitrary color missing at \(v\) under \(\pi\).
\end{lemma}
\begin{proof}
First suppose that no positive slot \(h\leq i\) is a PE supported at
\(v\).  The palette inequality and
\noderef{MissingColorsNonemptyOfMaximumDegree}, applied to the arbitrary
partial coloring \(\pi\) at \(v\), give
\(\gamma\in\overline\pi(v)\).  Together with the assumed absence of a
prior supporting PE, these are exactly the second alternative of
\noderef{ParallelExtensionColor}.

Now suppose at least one such PE exists.  Then \(i>0\); write \(q=i-1\).
Form the nonempty finite set
\[
 B=\{b\leq q:\ \Theta_{b+1}=\mathrm{PE}
                 \text{ and the support at slot }b+1\text{ is }v\}.
\]
List it in increasing order as
\(b_0<\cdots<b_p\).  Finiteness comes from the finite index type.
The exact-range property of this list says both that every listed slot
is a PE at \(v\), and that every PE at \(v\) through \(i\) occurs in the
list.

For each listed index, eliminate the certificate supplied by
\noderef{TashkinovSeries}.  Its kind equality excludes the initial, SE,
and RE constructors of \noderef{TashkinovSeriesStep}.  The remaining PE
constructor supplies exact colors \(\gamma_{b_a+1}\) and
\(\delta_{b_a+1}\), says that \(v\) lies in the tree processed at that
slot, and says that the resulting tree is an ordered extension.  Repeated
use of \noderef{TreeSequenceInitialSegment} therefore places \(v\) in
\(V(T_i)\), as required for the history-transport results.

Successive entries in the list obey
\[
             \gamma_{b_{a+1}+1}=\delta_{b_a+1}
\]
by \noderef{SupportingPEColorChain}.  More precisely, apply
\noderef{SupportingPEMissingTransport} to the entire exact enumeration,
with the first gamma \(\gamma_{b_0+1}\) and the last delta
\(\delta_{b_p+1}\).  Its second asserted equality is
\[
 \overline{\varphi}_{b_p+1}(v)
  =\bigl(\overline{\varphi}_{b_0}(v)
        \setminus\{\gamma_{b_0+1}\}\bigr)
       \cup\{\delta_{b_p+1}\}.                                 \tag{1}
\]
In particular, the last delta color is missing at \(v\) immediately
after the latest supporting PE.

Put \(h=b_p+1\).  By construction \(h\) is positive, \(h\leq i\),
\(\Theta_h=\mathrm{PE}\), its chosen support is \(v\), and
\(\delta_h=\delta_{b_p+1}\).  Every positive PE slot \(j\leq i\)
supported at \(v\) corresponds to an element of \(B\), hence its
zero-based position is at most the last element \(b_p\).  Thus \(j\leq h\);
this is exactly the required latest-slot clause.

It remains to show that \(\delta_h\) is missing under the selected
coloring \(\pi\), rather than only just after slot \(h\).  If \(h=i\),
equation (1) already gives missingness under \(\varphi_i\).  If \(h<i\),
the ordered-extension field at \(h\) places \(v\) in \(T_{h+1}\), and
maximality of \(h\) says that no PE in the interval
\(h+1,\ldots,i\) is supported at \(v\).  Apply
\noderef{HkuUnsupportedVertexColors} to this exact interval.  It yields
\[
          \overline{\varphi}_{h}(v)
             =\overline{\varphi}_{i}(v),                       \tag{2}
\]
so (1) again puts \(\delta_h\) in the right-hand set.  Finally the
selected-stability field of \noderef{MaximumDefectiveChoice}, using the
vertex clause of \noderef{StableColoring} at
\(v\in V(T_{i+1})\), gives
\[
                 \overline\pi(v)=\overline{\varphi}_i(v).       \tag{3}
\]
Equations (1)--(3) show
\(\delta_h\in\overline\pi(v)\).  Choose
\(\gamma=\delta_h\).  The positivity, index bound, kind, support, exact
delta equality, and latest-slot property proved above give the first
history alternative in \noderef{ParallelExtensionColor}; (3) gives its
missingness clause.  This completes both exhaustive cases.
\end{proof}
